% !TEX root = TA.tex
\chapter{HASIL LENGKAP UJI STATISTIK DOMAIN GENDER}
\label{chap:lampiran-statistik-gender}

Lampiran ini memuat seluruh nilai-p Exact McNemar mentah dan nilai-p setelah
koreksi Holm pada domain gender GSM8K, serta hasil tambahan ETHICS domain
gender. Sel berwarna menunjukkan nilai di bawah 0,05. Warna tersebut digunakan
untuk memudahkan pembacaan, sedangkan keputusan akhir tetap didasarkan pada
nilai setelah koreksi Holm dalam keluarga uji yang telah ditetapkan.

Pada gambar berikut, E merujuk pada persona eksplisit, I1--I3 merujuk pada
ketiga persona implisit, N merujuk pada skenario netral, L merujuk pada
laki-laki, dan P merujuk pada perempuan.

\begin{landscape}
\begin{figure}[H]
    \centering
    \includegraphics[
        width=\linewidth,
        height=0.82\textheight,
        keepaspectratio
    ]{image/bab6/gsm8k/gender/akurasi/17_gender_exact_mcnemar_and_holm_pvalue_heatmap.png}
    \caption[Nilai-p uji Exact McNemar pada domain gender]
    {Nilai-p uji Exact McNemar sebelum dan setelah koreksi Holm pada domain
    gender}
    \label{fig:lampiran-pvalue-gender}
\end{figure}
\end{landscape}

\section{Penggunaan Token Domain Gender}
\label{sec:lampiran-token-gender}

Gambar berikut menunjukkan rata-rata \textit{input token}, \textit{completion
token}, dan total token pada setiap model. M dan W masing-masing menunjukkan
persona laki-laki dan perempuan, E menunjukkan skenario eksplisit, I Avg
menunjukkan rata-rata I1--I3, dan N menunjukkan skenario netral.

\begin{figure}[H]
    \centering
    \begin{subfigure}{\linewidth}
        \centering
        \includegraphics[
            width=0.98\linewidth,
            height=0.22\textheight,
            keepaspectratio
        ]{image/BabVI/gsm8k/gender/gsm8k_gender_token_components_v5.png}
        \caption{Rata-rata komponen penggunaan token menurut model dan skenario
        persona}
        \label{fig:lampiran-token-gender-components}
    \end{subfigure}

    \vspace{0.25cm}

    \begin{subfigure}{\linewidth}
        \centering
        \includegraphics[
            width=0.98\linewidth,
            height=0.39\textheight,
            keepaspectratio
        ]{image/bab6/gsm8k/gender/akurasi/22_gender_completion_token_accuracy_tradeoff_implicit_vs_explicit.png}
        \caption{Perubahan \textit{completion token} dan akurasi skenario
        implisit terhadap skenario eksplisit}
        \label{fig:lampiran-token-akurasi-gender}
    \end{subfigure}

    \caption[Penggunaan token dan akurasi pada domain gender]
    {Penggunaan token pada domain gender dan perbandingannya dengan akurasi}
    \label{fig:lampiran-token-gender-gabungan}
\end{figure}

\section{Penanda Linguistik Domain Gender}
\label{sec:lampiran-marker-gender}

Bagian ini memuat hasil pendukung penanda linguistik pada domain gender.
Tabel dan gambar berikut digunakan untuk melengkapi pembahasan ringkas pada
Bab~\ref{chap:evaluasi}, terutama untuk penanda numerik, cara model menyusun
jawaban, jumlah teks yang dapat dibandingkan, dan proporsi kategori awal jawaban
pada perbandingan antarpersona.

Tabel~\ref{tab:lampiran-penanda-numerik-gender} memuat subpenanda numerik
dengan nilai mutlak \textit{effect size} \textit{rank-biserial} terbesar pada
setiap kelompok dan model. Kolom E--N sampai I3--N membandingkan setiap
skenario dengan skenario netral. Kolom laki-laki--perempuan mengambil nilai
mutlak terbesar dari perbandingan antarpersona pada E, I1, I2, atau I3. Tabel ini
menjadi hasil pendukung untuk pembahasan penanda linguistik pada
Bab~\ref{chap:evaluasi}.

\begin{landscape}
\scriptsize
\setlength{\LTpre}{0pt}
\setlength{\LTpost}{0pt}
\begin{longtable}{@{}>{\raggedright\arraybackslash}p{2.45cm}>{\raggedright\arraybackslash}p{2.1cm}>{\raggedright\arraybackslash}p{2.55cm}>{\raggedright\arraybackslash}p{2.55cm}>{\raggedright\arraybackslash}p{2.55cm}>{\raggedright\arraybackslash}p{2.55cm}>{\raggedright\arraybackslash}p{2.55cm}@{}}
\caption[Subpenanda numerik utama pada domain gender]
{Subpenanda numerik dengan ukuran efek terbesar pada domain gender}
\label{tab:lampiran-penanda-numerik-gender} \\
\toprule
Model & Kelompok penanda & E--N & I1--N & I2--N & I3--N & Laki-laki--perempuan \\
\midrule
\endfirsthead
\multicolumn{7}{c}{\tablename\ \thetable\ (lanjutan)} \\
\toprule
Model & Kelompok penanda & E--N & I1--N & I2--N & I3--N & Laki-laki--perempuan \\
\midrule
\endhead
\midrule
\multicolumn{7}{r}{\textit{Bersambung ke halaman berikutnya}} \\
\endfoot
\bottomrule
\endlastfoot
DeepSeek R1-Distill 32B & Panjang jawaban & panjang kalimat (r=-0.15) & jumlah kata (r=-0.37) & jumlah kata (r=-0.41) & jumlah kata (r=-0.45) & panjang kalimat (r=0.18) \\
DeepSeek R1-Distill 32B & Variasi kata & variasi kata (r=0.09) & variasi kata (r=0.48) & variasi kata (r=0.53) & variasi kata (r=0.54) & variasi kata (r=0.05) \\
DeepSeek R1-Distill 32B & Cara menyusun jawaban & langkah eksplisit (r=-0.29) & langkah eksplisit (r=-0.9) & langkah eksplisit (r=-0.91) & langkah eksplisit (r=-0.92) & langkah eksplisit (r=-0.2) \\
DeepSeek R1-Distill 32B & Penjelasan tambahan & penjelasan tambahan (r=-0.33) & penjelasan tambahan (r=0.47) & penjelasan tambahan (r=0.4) & penjelasan tambahan (r=0.4) & penjelasan tambahan (r=-0.4) \\
DeepSeek R1-Distill 32B & Bahasa matematis & aritmetika dalam kalimat (r=-0.08) & LaTeX \textit{display} (r=-1) & LaTeX \textit{display} (r=-0.99) & LaTeX \textit{display} (r=-0.99) & LaTeX \textit{display} (r=-0.24) \\
DeepSeek R1-Distill 32B & Kepastian dan keraguan & asumsi (r=1) & koreksi diri (r=1) & koreksi diri (r=1) & koreksi diri (r=1) & penegas (r=1) \\
DeepSeek R1-Distill 32B & Kata penghubung & kondisi (r=0.49) & sebab-akibat (r=0.65) & sebab-akibat (r=0.62) & sebab-akibat (r=0.58) & kontras (r=0.49) \\
DeepSeek R1-Distill 32B & Kata ganti diri & kata ganti orang pertama (r=1) & kata ganti orang pertama (r=1) & kata ganti orang pertama (r=1) & kata ganti orang pertama (r=1) & kata ganti orang pertama (r=-1) \\
DeepSeek R1-Distill 32B & Sapaan kepada pengguna & sapaan/ajakan (r=1) & sapaan/ajakan (r=1) & sapaan/ajakan (r=0) & sapaan/ajakan (r=1) & sapaan/ajakan (r=1) \\
Llama 3.1 8B & Panjang jawaban & panjang kalimat (r=0.08) & panjang kalimat (r=-0.29) & panjang kalimat (r=0.03) & panjang kalimat (r=-0.3) & panjang kalimat (r=-0.36) \\
Llama 3.1 8B & Variasi kata & variasi kata (r=0.03) & variasi kata (r=-0.04) & variasi kata (r=0.1) & variasi kata (r=0.07) & variasi kata (r=-0.09) \\
Llama 3.1 8B & Cara menyusun jawaban & subjudul (r=-0.5) & langkah eksplisit (r=0.46) & subjudul (r=-0.86) & langkah eksplisit (r=0.48) & subjudul (r=0.91) \\
Llama 3.1 8B & Penjelasan tambahan & penjelasan tambahan (r=0.62) & penjelasan tambahan (r=0.57) & penjelasan tambahan (r=0.57) & penjelasan tambahan (r=0.52) & penjelasan tambahan (r=0.57) \\
Llama 3.1 8B & Bahasa matematis & aritmetika dalam kalimat (r=-0.02) & aritmetika dalam kalimat (r=-0.13) & aritmetika dalam kalimat (r=-0.03) & aritmetika dalam kalimat (r=-0.11) & aritmetika dalam kalimat (r=-0.14) \\
Llama 3.1 8B & Kepastian dan keraguan & penegas (r=1) & penegas (r=1) & koreksi diri (r=0.33) & penegas (r=1) & keraguan (r=-1) \\
Llama 3.1 8B & Kata penghubung & sebab-akibat (r=-0.28) & kontras (r=0.1) & kontras (r=0.22) & kontras (r=0.21) & kondisi (r=0.3) \\
Llama 3.1 8B & Kata ganti diri & kata ganti orang pertama (r=0.64) & kata ganti orang pertama (r=-0.33) & kata ganti orang pertama (r=-1) & kata ganti orang pertama (r=-0.33) & kata ganti orang pertama (r=0.33) \\
Llama 3.1 8B & Sapaan kepada pengguna & sapaan/ajakan (r=-1) & sapaan/ajakan (r=-1) & sapaan/ajakan (r=-1) & sapaan/ajakan (r=-1) & sapaan/ajakan (r=-1) \\
Llama 3.3 70B & Panjang jawaban & panjang kalimat (r=-0.08) & panjang kalimat (r=-0.3) & jumlah kata (r=0.17) & panjang kalimat (r=-0.45) & jumlah kata (r=0.39) \\
Llama 3.3 70B & Variasi kata & variasi kata (r=0.01) & variasi kata (r=0.06) & variasi kata (r=0.09) & variasi kata (r=0.1) & variasi kata (r=-0.19) \\
Llama 3.3 70B & Cara menyusun jawaban & subjudul (r=-0.83) & daftar berpoin (r=-0.3) & subjudul (r=-0.94) & subjudul (r=-0.68) & subjudul (r=-1) \\
Llama 3.3 70B & Penjelasan tambahan & penjelasan tambahan (r=0.43) & penjelasan tambahan (r=0.58) & penjelasan tambahan (r=0.52) & penjelasan tambahan (r=0.58) & penjelasan tambahan (r=0.42) \\
Llama 3.3 70B & Bahasa matematis & LaTeX \textit{display} (r=-0.4) & LaTeX \textit{display} (r=-1) & LaTeX \textit{display} (r=-1) & LaTeX \textit{display} (r=-1) & LaTeX \textit{display} (r=1) \\
Llama 3.3 70B & Kepastian dan keraguan & penegas (r=-1) & penegas (r=0.52) & asumsi (r=0.35) & penegas (r=0.67) & koreksi diri (r=-1) \\
Llama 3.3 70B & Kata penghubung & kondisi (r=0.19) & kontras (r=0.3) & sebab-akibat (r=0.42) & sebab-akibat (r=0.32) & kondisi (r=-0.23) \\
Llama 3.3 70B & Kata ganti diri & kata ganti orang pertama (r=1) & kata ganti we (r=0.33) & kata ganti orang pertama (r=1) & kata ganti orang pertama (r=1) & kata ganti orang pertama (r=-1) \\
Llama 3.3 70B & Sapaan kepada pengguna & sapaan/ajakan (r=0.33) & sapaan/ajakan (r=0.6) & sapaan/ajakan (r=-1) & sapaan/ajakan (r=0.67) & sapaan/ajakan (r=1) \\
Ministral 8B & Panjang jawaban & jumlah kata (r=0.02) & panjang kalimat (r=0.14) & panjang kalimat (r=0.2) & panjang kalimat (r=0.18) & jumlah kata (r=-0.48) \\
Ministral 8B & Variasi kata & variasi kata (r=0.03) & variasi kata (r=-0.17) & variasi kata (r=-0.1) & variasi kata (r=-0.13) & variasi kata (r=0.41) \\
Ministral 8B & Cara menyusun jawaban & daftar berpoin (r=0.11) & langkah eksplisit (r=0.33) & langkah eksplisit (r=0.24) & langkah eksplisit (r=0.25) & subjudul (r=0.29) \\
Ministral 8B & Penjelasan tambahan & penjelasan tambahan (r=-0.12) & penjelasan tambahan (r=-0.2) & penjelasan tambahan (r=-0.27) & penjelasan tambahan (r=-0.25) & penjelasan tambahan (r=0.42) \\
Ministral 8B & Bahasa matematis & LaTeX \textit{display} (r=-0.18) & LaTeX \textit{display} (r=-0.15) & LaTeX \textit{display} (r=-0.16) & LaTeX \textit{display} (r=-0.09) & LaTeX \textit{display} (r=0.09) \\
Ministral 8B & Kepastian dan keraguan & asumsi (r=0.35) & penegas (r=0.34) & keraguan (r=-0.21) & koreksi diri (r=-0.28) & koreksi diri (r=-0.39) \\
Ministral 8B & Kata penghubung & kontras (r=0.15) & kontras (r=-0.29) & kontras (r=-0.35) & kontras (r=-0.29) & kontras (r=0.29) \\
Ministral 8B & Kata ganti diri & kata ganti orang pertama (r=-0.55) & kata ganti orang pertama (r=-0.82) & kata ganti orang pertama (r=-0.53) & kata ganti orang pertama (r=-0.6) & kata ganti orang pertama (r=-0.54) \\
Ministral 8B & Sapaan kepada pengguna & sapaan/ajakan (r=-1) & sapaan/ajakan (r=0.2) & sapaan/ajakan (r=-0.2) & sapaan/ajakan (r=-1) & sapaan/ajakan (r=-1) \\
GPT-OSS 20B & Panjang jawaban & jumlah kata (r=-0.37) & panjang kalimat (r=-0.28) & panjang kalimat (r=-0.27) & panjang kalimat (r=-0.3) & panjang kalimat (r=0.21) \\
GPT-OSS 20B & Variasi kata & variasi kata (r=0.35) & variasi kata (r=0.11) & variasi kata (r=-0.04) & variasi kata (r=0.13) & variasi kata (r=0.18) \\
GPT-OSS 20B & Cara menyusun jawaban & subjudul (r=0.38) & subjudul (r=0.81) & subjudul (r=0.9) & subjudul (r=0.86) & subjudul (r=-0.35) \\
GPT-OSS 20B & Penjelasan tambahan & penjelasan tambahan (r=-0.22) & penjelasan tambahan (r=-0.57) & penjelasan tambahan (r=-0.46) & penjelasan tambahan (r=-0.41) & penjelasan tambahan (r=0.24) \\
GPT-OSS 20B & Bahasa matematis & LaTeX \textit{display} (r=-0.46) & LaTeX \textit{display} (r=-0.14) & LaTeX \textit{display} (r=-0.14) & aritmetika dalam kalimat (r=-0.2) & aritmetika dalam kalimat (r=-0.13) \\
GPT-OSS 20B & Kepastian dan keraguan & penegas (r=-1) & penegas (r=-1) & koreksi diri (r=-0.45) & penegas (r=-1) & penegas (r=-1) \\
GPT-OSS 20B & Kata penghubung & kontras (r=-1) & kontras (r=1) & kontras (r=0.33) & kontras (r=0.33) & kontras (r=1) \\
GPT-OSS 20B & Kata ganti diri & kata ganti we (r=-0.46) & kata ganti we (r=-0.29) & kata ganti we (r=-0.12) & kata ganti we (r=-0.1) & kata ganti we (r=-0.29) \\
GPT-OSS 20B & Sapaan kepada pengguna & sapaan/ajakan (r=0) & sapaan/ajakan (r=0) & sapaan/ajakan (r=0) & sapaan/ajakan (r=0) & sapaan/ajakan (r=0) \\
Qwen3 Next 80B-A3B & Panjang jawaban & jumlah kata (r=-0.18) & jumlah kata (r=-0.2) & panjang kalimat (r=-0.15) & panjang kalimat (r=-0.17) & jumlah kata (r=-0.13) \\
Qwen3 Next 80B-A3B & Variasi kata & variasi kata (r=0.14) & variasi kata (r=0.17) & variasi kata (r=0.14) & variasi kata (r=0.14) & variasi kata (r=-0.07) \\
Qwen3 Next 80B-A3B & Cara menyusun jawaban & subjudul (r=-0.27) & subjudul (r=-0.6) & subjudul (r=-0.36) & subjudul (r=-0.36) & subjudul (r=-0.41) \\
Qwen3 Next 80B-A3B & Penjelasan tambahan & penjelasan tambahan (r=-0.08) & penjelasan tambahan (r=0.08) & penjelasan tambahan (r=0.06) & penjelasan tambahan (r=-0.13) & penjelasan tambahan (r=0.13) \\
Qwen3 Next 80B-A3B & Bahasa matematis & LaTeX \textit{display} (r=-0.15) & LaTeX \textit{display} (r=-0.27) & LaTeX \textit{display} (r=-0.21) & LaTeX \textit{display} (r=-0.23) & LaTeX \textit{display} (r=-0.09) \\
Qwen3 Next 80B-A3B & Kepastian dan keraguan & penegas (r=0.49) & penegas (r=0.58) & penegas (r=0.64) & penegas (r=-0.6) & penegas (r=0.78) \\
Qwen3 Next 80B-A3B & Kata penghubung & kontras (r=-0.15) & sebab-akibat (r=-0.26) & sebab-akibat (r=-0.26) & sebab-akibat (r=-0.29) & kontras (r=0.37) \\
Qwen3 Next 80B-A3B & Kata ganti diri & kata ganti orang pertama (r=0.23) & kata ganti orang pertama (r=0.43) & kata ganti orang pertama (r=0.47) & kata ganti orang pertama (r=0.39) & kata ganti orang pertama (r=0.31) \\
Qwen3 Next 80B-A3B & Sapaan kepada pengguna & sapaan/ajakan (r=0.24) & sapaan/ajakan (r=-0.67) & sapaan/ajakan (r=-0.67) & sapaan/ajakan (r=0.2) & sapaan/ajakan (r=-0.6) \\
\end{longtable}
\end{landscape}

Nilai $r$ yang besar pada penanda yang jarang muncul, seperti penjelasan
tambahan, dapat berasal dari sedikit pasangan yang berubah. Karena itu,
besarnya $r$ pada tabel ini digunakan sebagai hasil pendukung, bukan sebagai
satu-satunya dasar pembahasan. Pernyataan identitas AI, pola pembuka, dan
pronomina gender tidak dimasukkan ke tabel ini karena ketiganya dianalisis
sebagai kehadiran atau kategori, bukan sebagai densitas numerik dengan
\textit{rank-biserial}.

\begin{figure}[H]
    \centering
    \includegraphics[width=\textwidth]{image/BabVI/gsm8k/gender/consistent_20261002/03_numeric_effect_volcano_vs_neutral.png}
    \caption{Ukuran efek penanda linguistik numerik terhadap skenario netral pada domain gender}
    \label{fig:lampiran-marker-volcano-netral-gender}
\end{figure}

\begin{figure}[H]
    \centering
    \includegraphics[width=\textwidth]{image/BabVI/gsm8k/gender/consistent_20261002/15_direct_gender_effect_consistency.png}
    \caption[Ukuran efek penanda linguistik antarpersona gender]
    {Sebaran ukuran efek penanda linguistik numerik pada perbandingan persona
    laki-laki dan perempuan dalam domain gender}
    \label{fig:lampiran-marker-gender-laki-perempuan}
\end{figure}

Pada grafik langsung, warna membedakan model; segitiga menunjukkan
$p_{\mathrm{Holm}}<0{,}05$, sedangkan lingkaran menunjukkan hasil lainnya.
Ukuran titik mengikuti jumlah pasangan dengan selisih bukan nol dan berlian
menunjukkan median enam model. Arah ukuran efek tetap laki-laki dikurangi perempuan.


\begin{figure}[H]
    \centering
    \includegraphics[width=\textwidth]{image/BabVI/gsm8k/gender/consistent_20261002/05_proxy_presence_forest_structure_logic.png}
    \caption{Selisih penanda struktur dan logika antara skenario implisit dan netral}
    \label{fig:lampiran-marker-struktur-gender}
\end{figure}
\clearpage
\begin{figure}[H]
    \centering
    \includegraphics[width=\textwidth]{image/BabVI/gsm8k/gender/consistent_20261002/05_proxy_presence_forest_structure_logic_p02.png}
    \caption{Penghubung sebab-akibat: lanjutan perbandingan penanda linguistik gender}
    \label{fig:lampiran-marker-struktur-gender-p2}
\end{figure}

\clearpage
\begin{figure}[H]
    \centering
    \includegraphics[width=\textwidth]{image/BabVI/gsm8k/gender/consistent_20261002/05_proxy_presence_forest_structure_logic_p03.png}
    \caption{LaTeX \textit{display}: lanjutan perbandingan penanda linguistik gender}
    \label{fig:lampiran-marker-struktur-gender-p3}
\end{figure}


\clearpage

\begin{figure}[H]
    \centering
    \includegraphics[width=\textwidth]{image/BabVI/gsm8k/gender/consistent_20261002/16_direct_gender_availability_forest.png}
    \caption{Jumlah respons yang dapat dibandingkan antara persona laki-laki dan perempuan}
    \label{fig:lampiran-jumlah-respons-langsung-gender}
\end{figure}
\clearpage
\begin{figure}[H]
    \centering
    \includegraphics[width=\textwidth]{image/BabVI/gsm8k/gender/consistent_20261002/16_direct_gender_availability_forest_p02.png}
    \caption{Hanya konten teknis: lanjutan perbandingan penanda linguistik gender}
    \label{fig:lampiran-jumlah-respons-langsung-gender-p2}
\end{figure}


\begin{figure}[H]
    \centering
    \includegraphics[width=\textwidth]{image/BabVI/gsm8k/gender/consistent_20261002/17_direct_gender_opening_composition.png}
    \caption{Selisih kategori pembuka respons antara persona laki-laki dan perempuan}
    \label{fig:lampiran-pembuka-langsung-gender}
\end{figure}

\section{Penanda Linguistik ETHICS Domain Gender}
\label{sec:lampiran-ethics-gender-lm}

Gambar berikut melengkapi penanda linguistik ETHICS pada Bab VI. Seluruh
hasil masih berstatus \textit{candidate} dan belum menggantikan audit manusia
independen. Koreksi Holm mengikuti keluarga pengujian pada hasil sumber,
bukan dihitung ulang berdasarkan pembagian gambar. Pada grafik langsung,
segitiga menunjukkan $p_{\mathrm{Holm}}<0{,}05$ dan lingkaran menunjukkan
hasil lainnya; warna membedakan model.

\clearpage
\begin{figure}[H]
    \centering
    \includegraphics[width=\textwidth]{image/BabVI/ethics/gender/linguistic_marker_20261002/00_cakupan.png}
    \caption{Ketersediaan penjelasan dan pasangan analisis; L: laki-laki, P: perempuan (\textit{candidate})}
    \label{fig:ethics-gender-lm-cakupan}
\end{figure}

\clearpage
\begin{figure}[H]
    \centering
    \includegraphics[width=\textwidth]{image/BabVI/ethics/gender/linguistic_marker_20261002/02_panjang_penjelasan.png}
    \caption{Panjang penjelasan menurut skenario; titik: median, garis: rentang antarkuartil (\textit{candidate})}
    \label{fig:ethics-gender-lm-panjang}
\end{figure}

\clearpage
\begin{figure}[H]
    \centering
    \includegraphics[width=\textwidth]{image/BabVI/ethics/gender/linguistic_marker_20261002/N01_bahasa_umum_vs_netral.png}
    \caption{Ukuran efek penanda linguistik umum terhadap skenario netral (\textit{candidate})}
    \label{fig:ethics-gender-lm-umum-netral}
\end{figure}

\clearpage
\begin{figure}[H]
    \centering
    \includegraphics[width=\textwidth]{image/BabVI/ethics/gender/linguistic_marker_20261002/N03_wacana_gender_vs_netral.png}
    \caption{Ukuran efek penanda linguistik wacana gender terhadap skenario netral (\textit{candidate})}
    \label{fig:ethics-gender-lm-wacana-netral}
\end{figure}

\clearpage
\begin{figure}[H]
    \centering
    \includegraphics[width=\textwidth]{image/BabVI/ethics/gender/linguistic_marker_20261002/01_bahasa_umum_man_minus_woman.png}
    \caption{Perbedaan penanda linguistik umum antara persona laki-laki dan perempuan (\textit{candidate})}
    \label{fig:ethics-gender-lm-umum-langsung}
\end{figure}

\clearpage
\begin{figure}[H]
    \centering
    \includegraphics[width=\textwidth]{image/BabVI/ethics/gender/linguistic_marker_20261002/03_bahasa_gender_man_minus_woman.png}
    \caption{Perbedaan penanda linguistik wacana gender antara persona laki-laki dan perempuan (\textit{candidate})}
    \label{fig:ethics-gender-lm-wacana-langsung}
\end{figure}

\clearpage
\begin{figure}[H]
    \centering
    \includegraphics[width=\textwidth]{image/BabVI/ethics/gender/linguistic_marker_20261002/04_marker_biner_man_minus_woman.png}
    \caption{Selisih prevalensi penanda biner antara persona laki-laki dan perempuan; garis: interval kepercayaan 95\% marginal (\textit{candidate})}
    \label{fig:ethics-gender-lm-biner-langsung}
\end{figure}



\section{Hasil Tambahan ETHICS Domain Gender}
\label{sec:lampiran-ethics-gender}

\input{Lampiran-D-ETHICS-Akurasi-Lengkap.tex}

